\documentclass[11pt]{article}
\usepackage[a4paper,margin=1in]{geometry}
\usepackage{subfiles}
\usepackage{amsmath,amssymb,bm}
\usepackage{hyperref}
\usepackage{booktabs}
\usepackage{array}
\usepackage{mathtools}

\newcommand{\rs}{r_{\mathrm{s}}}

\title{Appendix 14 (to main theory): Bi-conformal Interior and the Space Deficit Interpretation}
\author{}
\date{}

\begin{document}
\let\phi\varphi
\maketitle

\section*{Purpose and scope}
This appendix develops the physical interpretation of the bi-conformal interior derived in the main theory (Sec.~VI--VIII). No new fields, potentials, or action modifications are introduced. Every statement follows from the same covariant action and bi-conformal ansatz used throughout the paper.

\textbf{Note on nomenclature.}  This appendix was
previously filed under the name ``ISPG\_Quantum''; it has
been renamed to ``ISPG\_Interior'' to avoid confusion with
the standalone quantum extension paper~\cite{Lotishvili2026Q},
which treats $\hbar$-dependent physics (canonical
quantization, loop corrections, decoherence, and mass
spectrum).  The content of the present appendix is
\emph{entirely classical}.

\section{The bi-conformal ``freezing'' effect}\label{sec:freezing}

The static bi-conformal metric with $\phi_0 = -\rs/r$ yields:
\begin{align}
g_{tt} &= -e^{-\rs/r}, &
g_{rr} &= e^{+\rs/r}.
\end{align}
As $r \to 0$, $e^{-\rs/r} \to 0$ exponentially and $e^{+\rs/r} \to \infty$ exponentially. This means:
\begin{itemize}
\item \textbf{Clock rate:} $d\tau/dt = \sqrt{|g_{tt}|} = e^{-\rs/(2r)} \to 0$. Local clocks stop exponentially.
\item \textbf{Radial distance:} $d\ell = \sqrt{g_{rr}}\,dr = e^{\rs/(2r)}\,dr \to \infty$. The proper radial distance to $r=0$ diverges---the centre is infinitely far away.
\item \textbf{Effective speed:} $c_{\rm eff}(r) = c\,e^{-\rs/r} \to 0$. All dynamical processes are exponentially suppressed.
\end{itemize}
This is the mechanism verified symbolically and numerically in Appendix~15 (Verification).

\section{Mass as integrated pressure deficit}\label{sec:mass}

In the main theory, the scalar field $\phi = -\rs/r$ encodes the local pressure state of the information substance (IS). The Schwarzschild radius $\rs = 2GM/c^2$ sets the depth of the pressure deficit. The total gravitational mass is therefore
\begin{equation}
M = \frac{c^2\,\rs}{2G},
\end{equation}
which is the integrated effect of the spatial pressure rarefaction. In the ISPG framework, mass \emph{is} the deficit: a deeper and more extended rarefaction corresponds to a larger gravitational mass.

\section{Light does not disappear}\label{sec:light}

A photon falling radially inward experiences the coordinate speed $dr/dt = c\,e^{-\rs/r}$, which approaches zero exponentially as $r \to 0$. From the perspective of a distant observer, the photon takes an infinite coordinate time to reach $r = 0$. In terms of affine parameter, however, the geodesic reaches $r = 0$ in a finite interval (see Appendix~12), so the spacetime is geodesically \emph{incomplete}. The crucial distinction is that the boundary at $r = 0$ is curvature-regular ($K \to 0$), not singular: the incompleteness reflects the topology of the bi-conformal manifold, not a breakdown of physics.

The observational consequence is that light entering the deep-deficit region undergoes extreme redshift: $1 + z = e^{\rs/(2r)} \to \infty$. The photon is not absorbed by a finite-radius horizon; it becomes increasingly undetectable to external observers. This removes the finite-radius event-horizon ingredient of the standard information-loss argument in the static exterior, but it does not by itself prove information recovery. The fate of information depends on the open boundary condition or extension at $r = 0$.

\section{Threshold for horizon-like appearance (heuristic)}\label{sec:minmass}

The bi-conformal metric does not produce a true event horizon in the finite-radius static exterior sense discussed in the main theory. There is no special soft-horizon surface at $r = \rs$; this radius is the photon sphere and has only the finite redshift $1+z=e^{1/2}$. Horizon-like invisibility, if used as an observational descriptor, is a threshold-dependent deep-interior effect at radii where $e^{\rs/(2r)}$ exceeds the chosen redshift threshold:
\begin{equation}
r_{\rm th}(z_{\rm th})=\frac{\rs}{2\ln(1+z_{\rm th})}.
\end{equation}

The present appendix does not derive a minimum mass or a collapse criterion for horizon-like appearance. At most, the exterior metric supplies a heuristic redshift threshold: weak fields ($\rs/r \ll 1$) are perturbatively close to Minkowski, while very deep interior radii can become observationally invisible for a chosen redshift threshold. A source model and collapse calculation are required before this can be promoted to a minimum-mass result.

\section{After dissipation: the permanent space deficit}\label{sec:permanent}

If the compact object loses energy (e.g.\ through radiative processes), the question arises: what remains?

In the bi-conformal framework, the answer follows from the freezing effect itself. As the deficit region $\phi \to -\infty$ deepens, the effective speed $c_{\rm eff} \to 0$ and the local relaxation timescale diverges:
\begin{equation}
\tau_{\rm relax}(r) \;\propto\; e^{+\rs/r} \;\to\; \infty \quad\text{as } r \to 0.
\end{equation}
The deep interior cannot refill because the refilling process itself is exponentially suppressed by the very deficit it would need to erase. This is a self-reinforcing mechanism: the deeper the deficit, the slower the recovery, and therefore the more permanent the deficit.

The endpoint is not a material relic of finite mass. It is a \textbf{permanent space deficit}---a region of arbitrarily deep rarefaction that persists indefinitely because its own depth prevents external processes from reaching it on any finite timescale.

\section{Compatibility with the main theory}\label{sec:compat}

This interpretation introduces no new physics. It follows directly from:
\begin{itemize}
\item the bi-conformal metric ansatz (main theory, Eq.~(1));
\item the vacuum field equation $\Box\phi = 0$ (main theory, Eq.~(6));
\item the singularity resolution (main theory, Sec.~VIII): $K \to 0$, curvature-regular boundary;
\item the stability verification (Appendix~15): $c_{\rm eff}^2 > 0$, no tachyonic instability.
\end{itemize}
The space deficit picture replaces the need for Planck-scale relics or evaporation endpoint arguments: the bi-conformal structure itself provides a singularity-free, curvature-regular, and permanently stable interior without invoking $\hbar$ or semiclassical physics.

\section*{Takeaway}
The bi-conformal interior is a region of exponentially suppressed dynamics: clocks stop, distances diverge, and effective speeds vanish. Mass is the integrated pressure deficit of the information substance. Light entering the deep interior is asymptotically redshifted rather than absorbed by a finite-radius horizon; actual recovery remains boundary-condition dependent. After energy dissipation, the remaining space deficit is self-perpetuating because the refilling timescale diverges with the deficit depth. No new fields, no semiclassical arguments, and no material relics are introduced in this classical ISPG interpretation.
\end{document}
